\chapter{Bütünleştirici Veri Analizi Laboratuvarı}
\index{veri analizi iş akışı}
\index{yeniden üretilebilir analiz}

\begin{hedefler}
Bu bölümün sonunda okuyucu:
\begin{itemize}
  \item araştırma sorularını hedef değişken, karşılaştırma ve uygun yöntemle eşleştirebilecek,
  \item analizden önce veri sözlüğü, aralık, eksik değer, yinelenen kayıt ve tutarlılık kontrollerini yapabilecek,
  \item betimsel istatistik, güven aralığı, eşleştirilmiş ve bağımsız karşılaştırma, korelasyon ve regresyonu tek bir projede birleştirebilecek,
  \item varsayım incelemesini test sonucundan önce ve araştırma tasarımıyla birlikte yürütebilecek,
  \item aynı veri üzerinde farklı yöntemlerin neden farklı soruları yanıtlayabileceğini açıklayabilecek,
  \item etki büyüklüğü, güven aralığı ve $p$-değerini dengeli biçimde yorumlayabilecek,
  \item Python ve SPSS ile yeniden üretilebilir bir analiz akışı kurabilecek,
  \item bulguları tablo, grafik, yöntem paragrafı ve ölçülü sonuç diliyle bütünleştirebilecektir.
\end{itemize}
\end{hedefler}

\begin{hazirlik}
Bir veri dosyasını açtıktan sonra hipotez testine geçmeden önce yapılması
gereken en az altı kontrolü listeleyin. Her kontrol atlanırsa oluşabilecek bir
yanlış sonuca örnek verin.
\end{hazirlik}

\section{Laboratuvarın amacı}

Önceki bölümlerde yöntemler ayrı başlıklar altında incelendi. Gerçek araştırma
ise bir menüden test seçme işlemi değildir. Araştırma sorusunun tanımlanması,
verinin yapısının anlaşılması, temizleme kararlarının kaydedilmesi, varsayım
incelemesi, tahmin, görselleştirme ve raporlama birbirine bağlıdır.
Uygulamalı istatistikte bu bütünlüklü proje yaklaşımı, tek tek yöntemlerin
mekanik kullanımından daha belirleyicidir \cite{mccullagh2022,james2021}.

Bu laboratuvarın temel ilkesi şudur:
\[
\text{Soru}\longrightarrow\text{Tasarım}\longrightarrow\text{Veri}
\longrightarrow\text{Yöntem}\longrightarrow\text{Yorum}.
\]

Akışın yönü sonuç görüldükten sonra geriye çevrilmemelidir. Önce düşük
$p$-değeri bulunup sonra buna uygun araştırma sorusu yazmak doğrulayıcı analiz
değildir.

\begin{figure}[htbp]
\centering
\resizebox{0.96\linewidth}{!}{%
\begin{tikzpicture}[alt={Araştırma sorusundan tasarım, veri kalitesi, analiz ve raporlamaya uzanan; yeni araştırmaya geri bildirim veren veri analizi akışı},
  box/.style={draw=PrismBlue,fill=PrismLight,rounded corners,minimum width=2.4cm,minimum height=0.85cm,align=center,font=\small},
  arr/.style={->,thick,PrismBlue}
]
\node[box] (q) at (0,0) {Araştırma\\sorusu};
\node[box] (d) at (3.1,0) {Tasarım ve\\örnekleme};
\node[box] (c) at (6.2,0) {Veri kalitesi\\kontrolü};
\node[box] (a) at (9.3,0) {Analiz ve\\tanılar};
\node[box] (r) at (12.4,0) {Raporlama ve\\sınırlılıklar};
\draw[arr] (q)--(d);
\draw[arr] (d)--(c);
\draw[arr] (c)--(a);
\draw[arr] (a)--(r);
\draw[->,dashed,PrismSubheadingBlue] (r.south) to[bend left=28]
  node[below,font=\scriptsize,text=PrismGray]{şeffaf geri bildirim ve yeni araştırma} (q.south);
\end{tikzpicture}%
}
\caption{Bütünleştirici ve tekrarlanabilir veri analizi akışı.}
\end{figure}

\section{Araştırma problemi}

Bir eğitim birimi, iki öğretim yaklaşımının öğrenci başarısıyla ilişkisini ve
haftalık çalışma süresinin son-test puanını ne ölçüde yordadığını incelemek
istemektedir. On altı öğrencinin uygulama öncesi ve sonrası başarı puanları ile
haftalık çalışma süreleri kaydedilmiştir. İlk sekiz öğrenci A, sonraki sekiz
öğrenci B grubundadır.

Bu küçük veri seti öğretim amacıyla tasarlanmıştır. Gerçek bir çalışmada
örnekleme yöntemi, grup ataması, ölçme aracının geçerlik ve güvenirliği, etik
izin, bilgilendirilmiş onam ve veri güvenliği ayrıca belgelenmelidir.

\section{Veri sözlüğü ve analiz birimi}
\index{veri sözlüğü}
\index{analiz birimi}

Her satır bir öğrenciyi temsil eder. Aynı öğrencinin ön-test ve son-test
puanları aynı satırda tutulduğu için eşleştirme yapısı korunur.

\begin{table}[htbp]
\centering
\caption{Laboratuvar veri setinin kısa veri sözlüğü.}
\begin{tabular}{@{}p{0.16\textwidth}p{0.27\textwidth}p{0.18\textwidth}p{0.24\textwidth}@{}}
\toprule
Değişken & Açıklama & Tür & Kabul edilen değer \\
\midrule
\texttt{id} & Öğrencinin anonim kodu & Kimlik & 1--16, benzersiz \\
\texttt{grup} & Öğretim yaklaşımı & Kategorik & A veya B \\
\texttt{on} & Uygulama öncesi başarı & Nicel & 0--100 \\
\texttt{son} & Uygulama sonrası başarı & Nicel & 0--100 \\
\texttt{saat} & Haftalık çalışma süresi & Nicel & 0--40 saat \\
\bottomrule
\end{tabular}
\end{table}

Veri sözlüğü yalnız değişken adlarını açıklamaz; olanaksız değerleri, eksik
kodlarını, ölçü birimini ve kategori referanslarını da tanımlar. Örneğin 999
değeri eksik veri kodu olarak kullanılmışsa sıradan puan gibi analiz edilmeden
önce dönüştürülmelidir.

\section{Ham veri}

\begin{table}[htbp]
\centering
\caption{Bütünleştirici laboratuvar veri seti.}
\small
\begin{tabular}{rrrrr@{\qquad}rrrrr}
\toprule
No & Grup & Ön & Son & Saat & No & Grup & Ön & Son & Saat \\
\midrule
1&A&52&60&5 & 9&B&54&57&4\\
2&A&48&54&3 &10&B&57&61&5\\
3&A&61&68&7 &11&B&49&52&3\\
4&A&55&63&6 &12&B&60&64&7\\
5&A&58&65&6 &13&B&51&55&3\\
6&A&46&52&2 &14&B&56&60&5\\
7&A&63&70&8 &15&B&47&50&2\\
8&A&50&58&4 &16&B&59&63&6\\
\bottomrule
\end{tabular}
\end{table}

\section{Analizden önce veri kalite planı}
\index{veri kalitesi}
\index{veri temizleme}
\index{yinelenen kayıt}
\index{olanaksız değer}

Ham dosya açıldığında önce şu kontroller yapılır:

\begin{enumerate}
  \item Sütun adları ve veri türleri veri sözlüğüyle uyumlu mu?
  \item Her öğrenci kodu benzersiz mi; yinelenen veya atlanan kayıt var mı?
  \item Grup kodlarında boşluk, küçük harf veya bilinmeyen kategori var mı?
  \item Ön-test ve son-test puanları 0--100 aralığında mı?
  \item Çalışma süresi gerçekçi aralıkta ve doğru birimde mi?
  \item Eksik değer kodları gerçek eksik değere dönüştürülmüş mü?
  \item Ön-test ve son-test aynı öğrenciyle doğru eşleşmiş mi?
  \item Aykırı görünen değer veri hatası mı, gerçek gözlem mi?
\end{enumerate}

Bu veri setinde 16 benzersiz öğrenci, sekizer kişilik iki grup ve eksiksiz
kayıt vardır. En küçük ve en büyük değerler veri sözlüğündeki sınırlar
içindedir. Bununla birlikte küçük örneklem ve yapay veri niteliği genelleme
gücünü sınırlar.

\begin{mistakebox}
Veri temizleme, sonuçları daha anlamlı yapmak için gözlem silme işlemi değildir.
Her düzeltme özgün kayıtla doğrulanmalı, kural sonuç görülmeden tanımlanmalı ve
analiz günlüğünde saklanmalıdır.
\end{mistakebox}

\section{Araştırma soruları ve yöntem eşleştirmesi}
\index{yöntem seçimi}
\index{hedef büyüklük}

\begin{table}[htbp]
\centering
\caption{Araştırma sorusu, hedef büyüklük ve analiz yöntemi.}
\begingroup\scriptsize
\begin{tabular}{@{}p{0.32\textwidth}p{0.23\textwidth}p{0.27\textwidth}@{}}
\toprule
Araştırma sorusu & Hedef büyüklük & Temel yöntem \\
\midrule
Son-test puanları nasıl dağılmaktadır? & Merkez, yayılım, biçim & Grafik ve betimsel istatistik \\
Ortalama son-test puanı ne kadar hassastır? & Evren ortalaması & $t$ güven aralığı \\
Öğrenciler uygulama sonrasında değişmiş midir? & Ortalama eşleştirilmiş fark & Eşleştirilmiş $t$ testi \\
A ve B son-test ortalamaları farklı mıdır? & Bağımsız ortalama farkı & Welch $t$ testi \\
Grupların ortalama değişimleri farklı mıdır? & Değişim farkı & Welch testi veya başlangıç ayarlı model \\
Saat ile son-test ilişkili midir? & Pearson korelasyonu & Saçılım grafiği ve korelasyon \\
Saat son-test puanını ne ölçüde yordamaktadır? & Regresyon eğimi & Basit doğrusal regresyon \\
\bottomrule
\end{tabular}
\endgroup
\end{table}

Her yöntem farklı bir hedefi yanıtlar. Eşleştirilmiş testte bütün öğrencilerin
ortalama değişimi, bağımsız testte iki grubun ortalama farkı, regresyonda ise
çalışma saatine göre koşullu ortalama incelenir.

\section{Analiz planını sonuçtan önce yazmak}
\index{analiz planı}
\index{ön kayıt}

Kısa bir analiz planı şu kararları önceden belirtmelidir:

\begin{itemize}
  \item birincil ve ikincil araştırma soruları,
  \item değişkenlerin kodlanması ve farkın yönü,
  \item dışlama ve eksik veri kuralları,
  \item temel grafik ve varsayım incelemeleri,
  \item kullanılacak test, güven düzeyi ve etki büyüklüğü,
  \item çoklu analizlerin keşfedici veya doğrulayıcı niteliği,
  \item duyarlılık analizlerinin amacı.
\end{itemize}

Bu laboratuvarda fark $D=\text{son}-\text{ön}$ olarak tanımlanacaktır. Pozitif
değer artışı gösterir. Bütün testler öğretim amacıyla çift yönlü ve yüzde 95
güven düzeyinde ele alınacaktır.

\begin{researchbox}
Aynı küçük veri setinde çok sayıda soru sınanmaktadır. Bu nedenle sonuçlar
bağımsız doğrulayıcı kanıtlar gibi sunulmamalıdır. Laboratuvarın amacı yöntem
bağlantılarını göstermek; gerçek araştırmada birincil sonucu önceden
belirlemek gerekir.
\end{researchbox}

\section{Betimsel analiz}

Son-test puanları için
\[
\bar x=59{,}50,
\qquad s=5{,}90,
\qquad Md=59{,}0,
\qquad \min=50,
\qquad \max=70
\]
bulunur. Ortalama ile medyan yakındır ve veri aralığında açıkça olanaksız bir
değer yoktur.

Grup bazında son-test özetleri:

\begin{table}[htbp]
\centering
\caption{Gruplara göre ön-test, son-test ve değişim özetleri.}
\begin{tabular}{@{}lrrrrrr@{}}
\toprule
 & \multicolumn{2}{c}{Ön-test} & \multicolumn{2}{c}{Son-test} & \multicolumn{2}{c}{Değişim} \\
\cmidrule(lr){2-3}\cmidrule(lr){4-5}\cmidrule(l){6-7}
Grup & Ort. & SS & Ort. & SS & Ort. & SS \\
\midrule
A & 54,13 & 6,36 & 61,25 & 6,43 & 7,13 & 0,83 \\
B & 54,13 & 4,58 & 57,75 & 5,12 & 3,63 & 0,52 \\
\bottomrule
\end{tabular}
\end{table}

İki grubun ön-test ortalamaları bu yapay örnekte aynıdır. A grubunun ortalama
artışı B grubundan daha büyüktür. Yalnız ortalamaya bakmak yeterli değildir;
ham noktalar, değişim dağılımı ve küçük örneklem belirsizliği de incelenmelidir.

\section{Görselleştirme planı}
\index{veri görselleştirme}
\index{bağlantı grafiği}

Bu veri için tek bir grafik bütün soruları yanıtlamaz:

\begin{itemize}
  \item son-test dağılımı için nokta veya histogram,
  \item gruplar için yan yana nokta ve kutu grafiği,
  \item bireysel değişim için ön--son bağlantı grafiği,
  \item çalışma saati ile son-test için saçılım grafiği ve doğrusal uyum,
  \item regresyon tanısı için artık--uydurulan değer grafiği.
\end{itemize}

Ön--son bağlantı grafiği ortalama artışın her öğrenci için aynı olup olmadığını
gösterir. Bu veride bütün öğrencilerin puanı artmıştır; fakat A grubundaki
artışlar 6--8, B grubundakiler 3--4 puan arasındadır.

\begin{etkinlik}
Aynı veri için yalnız grup ortalamalarını gösteren çubuk grafik ile bütün
öğrenci noktalarını gösteren grafiği karşılaştırın. Küçük örneklemde hangi
bilginin çubuk grafikte kaybolduğunu açıklayın.
\end{etkinlik}

\section{Ortalama için güven aralığı}

Son-test ortalamasının standart hatası
\[
SE(\bar X)=\frac{5{,}90}{\sqrt{16}}=1{,}475
\]
ve 15 serbestlik derecesinde kritik değer yaklaşık 2,131'dir. Yüzde 95 güven
aralığı
\[
59{,}50\pm2{,}131(1{,}475)
=[56{,}36;\ 62{,}64]
\]
olarak bulunur.

Bu aralık örnekleme belirsizliğini gösterir. Kolayda seçilmiş, yanlı veya
geçersiz ölçülmüş bir örneklemin sorunlarını kapsamaz. Ayrıca öğrencilerin
yüzde 95'inin bu aralıkta olduğu anlamına gelmez.

\section{Eşleştirilmiş ön-test--son-test analizi}
\index{t testi@$t$ testi!eşleştirilmiş örneklem}
\index{ön-test--son-test tasarımı}

Her öğrenci için son-testten ön-test çıkarılır:
\[
D_i=\text{Son}_i-\text{Ön}_i.
\]
Farkların ortalaması ve standart sapması
\[
\bar d=5{,}375,
\qquad s_D=1{,}928
\]
olup standart hata
\[
SE(\bar D)=\frac{1{,}928}{\sqrt{16}}=0{,}482
\]
bulunur. Test istatistiği
\[
t=\frac{5{,}375}{0{,}482}=11{,}15,
\qquad sd=15,
\qquad p<0{,}001
\]
ve yüzde 95 güven aralığı
\[
[4{,}35;\ 6{,}40]
\]
olur.

Eşleştirilmiş standartlaştırılmış değişim
\[
d_z=\frac{5{,}375}{1{,}928}=2{,}79
\]
olarak hesaplanır. Bu yapay veride fark değişkenliği çok küçük olduğu için
etki olağandışı büyüktür; gerçek araştırmada bu büyüklüğün genellenebilirliği
ayrı örneklemde değerlendirilmelidir.

\begin{mistakebox}
Ön-test--son-test farkının anlamlı olması artışın uygulamadan kaynaklandığını
kanıtlamaz. Zaman, sınanma etkisi, olgunlaşma ve dış olaylar da değişim
üretebilir. Nedensel yorum için uygun karşılaştırma grubu ve atama süreci
gerekir.
\end{mistakebox}

\section{Bağımsız grupların son-test karşılaştırması}
\index{Welch t testi@Welch $t$ testi}

A grubunun son-test ortalaması 61,25, B grubunun 57,75'tir. Fark
\[
61{,}25-57{,}75=3{,}50
\]
puandır. Welch standart hatası yaklaşık 2,91 ve serbestlik derecesi 13,33'tür:
\[
t(13{,}33)=1{,}20,
\qquad p=0{,}249.
\]
Fark için yüzde 95 güven aralığı
\[
[-2{,}76;\ 9{,}76]
\]
olur. Cohen'in birleşik standart sapmaya dayalı $d$ değeri yaklaşık 0,60'tır.

Nokta tahmini A grubunda daha yüksek ortalama ve orta büyüklükte
standartlaştırılmış fark göstermesine rağmen aralık sıfırı ve ters yöndeki
küçük farkları da içermektedir. ``Gruplar eşittir'' değil, bu küçük örneklemde
son-test farkının hassas tahmin edilemediği söylenmelidir.

\begin{researchbox}
$d=0{,}60$ ile $p=0{,}249$ çelişki değildir. Etki büyüklüğü gözlenen farkın
derecesini, $p$-değeri bu farkın standart hataya göre ne kadar olağandışı
olduğunu gösterir. Sekizer kişilik gruplarda belirsizlik geniştir.
\end{researchbox}

\section{Değişim farkı ve başlangıç ayarı}
\index{değişim puanı}
\index{başlangıç ayarlı model}
\index{ANCOVA}

Son-test karşılaştırması grupların dönem sonundaki farkını; değişim puanı
karşılaştırması grupların ortalama artış farkını hedefler. A grubunun ortalama
artışı 7,125, B grubunun 3,625 olduğundan değişim farkı 3,50 puandır.

Değişim puanlarının grup içi standart sapmaları küçük olduğu için Welch testi
\[
t(11{,}69)=10{,}08,
\qquad p<0{,}001
\]
sonucunu verir. Aynı 3,50 puanlık farkın son-test analizinde belirgin değil,
değişim analizinde belirgin olması paydadaki standart hatanın farklılığından
kaynaklanır.

Başka bir yaklaşım, son-test puanını ön-test ve grup ile modelleyen başlangıç
ayarlı regresyon veya ANCOVA'dır:
\[
\text{Son}=\beta_0+\beta_1\text{Ön}+\beta_2\text{Grup}+\varepsilon.
\]
Bu yapay veride ön-test son-testin çok güçlü yordayıcısıdır ve ayarlı grup
farkı yine 3,50 puandır. Ancak değişim analizi ile ANCOVA her veri yapısında
aynı sonucu vermez.

\begin{mistakebox}
Son-test, değişim puanı ve ANCOVA arasından en küçük $p$-değerini veren yöntemi
seçmeyin. Hedef büyüklük, randomizasyon, başlangıç dengesi, ölçüm güvenirliği
ve analiz planı yöntem seçiminden önce belirlenmelidir.
\end{mistakebox}

\section{Korelasyon analizi}
\index{Pearson korelasyonu}

Çalışma süresi ile son-test puanı arasında
\[
r(14)=0{,}981,
\qquad p<0{,}001
\]
bulunur. Saçılım grafiği güçlü ve yaklaşık doğrusal pozitif örüntü gösterir.
Basit doğrusal modelde $r^2=0{,}962$ olur.

Bu olağandışı yüksek ilişki küçük ve öğretim için oluşturulmuş veri setinin
özelliğidir. Gerçek çalışmada güven aralığı, etkili gözlemler, ölçüm
güvenirliği ve yeni örneklem doğrulaması özellikle önemlidir.

Korelasyon çalışma süresinin başarıya neden olduğunu kanıtlamaz. Başlangıç
başarısı, motivasyon, ders devamı veya aile desteği her iki değişkenle ilişkili
olabilir. Ayrıca başarılı öğrenciler daha fazla çalışma bildirebilir.

\section{Basit doğrusal regresyon}
\index{basit doğrusal regresyon}

Uydurulan doğru
\[
\widehat{\text{Son puan}}
=44{,}60+3{,}14\times\text{Saat}
\]
biçimindedir. Eğimin standart hatası 0,168 ve yüzde 95 güven aralığı
\[
[2{,}78;\ 3{,}50]
\]
olarak bulunur. Test sonucu
\[
t(14)=18{,}73,
\qquad p<0{,}001
\]
ve model uyumu
\[
R^2=0{,}962,
\qquad s_e=1{,}20
\]
şeklindedir.

Gözlenen 2--8 saat aralığında çalışma süresi bir saat daha yüksek olan
öğrencilerin beklenen son-test puanı yaklaşık 3,14 puan daha yüksektir. Bu
ilişki 15 saat için doğrudan uzatılmamalı ve nedensel ``bir saat çalıştırmak
3,14 puan artırır'' ifadesine dönüştürülmemelidir.

\section{Kategorik bir türetim: Bilgi kaybını görmek}
\index{kategorileştirme}
\index{bilgi kaybı}

Değişim puanı en az 7 olan öğrenciler ``yüksek artış'' olarak sınıflandırılsın.
A grubunda 6, B grubunda 0 öğrenci yüksek artış göstermiştir:

\begin{table}[htbp]
\centering
\caption{Grup ve yüksek artış sınıflandırması.}
\begin{tabular}{@{}lrrr@{}}
\toprule
 & Yüksek artış & Diğer & Toplam \\
\midrule
A & 6 & 2 & 8 \\
B & 0 & 8 & 8 \\
\bottomrule
\end{tabular}
\end{table}

Beklenen frekanslardan ikisi 5'in altındadır ve bir gözlenen hücre sıfırdır;
Fisher kesin testi uygundur ve iki yönlü $p=0{,}007$ verir. Bununla birlikte 7
puan eşiği veri görüldükten sonra seçilmişse sonuç doğrulayıcı değildir.

Nicel değişim puanını iki kategoriye ayırmak değerler arasındaki farkları
kaybeder. 7 ve 8 puan aynı, 6 ve 7 puan tamamen farklı kabul edilir. Klinik
veya eğitsel olarak doğrulanmış eşik yoksa özgün nicel analiz önceliklidir.

\section{Varsayım ve duyarlılık matrisi}
\index{varsayım kontrolü}
\index{duyarlılık analizi}

\begin{table}[htbp]
\centering
\caption{Laboratuvar analizlerinde temel tanı soruları.}
\begingroup\scriptsize
\begin{tabular}{@{}p{0.20\textwidth}p{0.28\textwidth}p{0.34\textwidth}@{}}
\toprule
Analiz & Temel kontrol & Sorun varsa olası yaklaşım \\
\midrule
Güven aralığı & Örnekleme, aykırı değer, dağılım & Bootstrap veya sağlam aralık \\
Eşleştirilmiş $t$ & Farkların bağımsızlığı ve biçimi & Wilcoxon, işaret veya sağlam yöntem \\
Welch $t$ & Bağımsızlık, aykırı değer, biçim & Mann--Whitney ya da sağlam ortalama \\
Korelasyon & Doğrusallık ve etkili gözlem & Spearman, eğrisel model, duyarlılık \\
Regresyon & Artık örüntüsü, varyans, etki & Sağlam SE, dönüşüm, farklı model \\
Kategorik test & Bağımsızlık ve beklenen frekans & Fisher veya kesin/Monte Carlo test \\
\bottomrule
\end{tabular}
\endgroup
\end{table}

Alternatif yöntem seçimi sonuç görüldükten sonra anlamlılık arama yolu
olmamalıdır. Ana analiz, varsayım sorunu ve duyarlılık analizinin amacı açıkça
ayrılmalıdır.

\section{Eksik veri senaryosu}
\index{eksik veri}
\index{seçici kayıp}

Bu öğretim verisi eksiksizdir; gerçek projede bu durum seyrektir. Örneğin B
grubundaki düşük son-test puanlı iki öğrenci izlem ölçümüne katılmasaydı grup
ortalaması yapay biçimde yükselebilirdi. Eksiklik yalnız örneklem hacmini
azaltmaz, sonuçla ilişkiliyse yanlılık da üretir.

Asgari raporlama:

\begin{itemize}
  \item her değişken için eksik sayı ve yüzde,
  \item gruplara göre kayıp örüntüsü,
  \item tam durum analizine giren öğrenci sayısı,
  \item kullanılan silme veya atama yöntemi,
  \item makul alternatiflere duyarlılık analizi.
\end{itemize}

Eksik değerleri sıfırla veya grup ortalamasıyla doldurmak değişkenliği ve
standart hatayı bozar. Yöntem eksiklik mekanizması ve araştırma tasarımına göre
seçilmelidir.

\section{Yeniden üretilebilir Python analizi}

\begin{lstlisting}
import numpy as np
import pandas as pd
from scipy import stats
import statsmodels.formula.api as smf

veri = pd.DataFrame({
    "id": range(1, 17),
    "grup": ["A"] * 8 + ["B"] * 8,
    "on": [52,48,61,55,58,46,63,50,
           54,57,49,60,51,56,47,59],
    "son": [60,54,68,63,65,52,70,58,
            57,61,52,64,55,60,50,63],
    "saat": [5,3,7,6,6,2,8,4,
             4,5,3,7,3,5,2,6]
})

# Veri kalite kontrolleri
assert veri["id"].is_unique
assert veri["grup"].isin(["A", "B"]).all()
assert veri[["on", "son"]].apply(
    lambda x: x.between(0, 100)).all().all()
assert veri["saat"].between(0, 40).all()
print(veri.isna().sum())

# Turetilmis degisken
veri["degisim"] = veri["son"] - veri["on"]
print(veri.groupby("grup")[["on", "son", "degisim"]]
      .agg(["count", "mean", "std"]))

# Guven araligi ve temel testler
ga = stats.t.interval(
    .95, len(veri) - 1,
    loc=veri["son"].mean(),
    scale=stats.sem(veri["son"])
)
esli = stats.ttest_rel(veri["son"], veri["on"])
A = veri.loc[veri["grup"] == "A", "son"]
B = veri.loc[veri["grup"] == "B", "son"]
welch = stats.ttest_ind(A, B, equal_var=False)
korelasyon = stats.pearsonr(veri["saat"], veri["son"])

# Regresyon ve baslangic ayarli model
basit = smf.ols("son ~ saat", data=veri).fit()
ayarli = smf.ols("son ~ on + C(grup)", data=veri).fit()

# Kesin kategorik analiz
veri["yuksek_artis"] = veri["degisim"] >= 7
tablo = pd.crosstab(veri["grup"], veri["yuksek_artis"])
fisher = stats.fisher_exact(tablo.to_numpy())

print("Yuzde 95 GA:", ga)
print("Eslestirilmis test:", esli)
print("Welch testi:", welch)
print("Korelasyon:", korelasyon)
print(basit.summary())
print(ayarli.summary())
print(tablo, fisher)
\end{lstlisting}

Kodda veri kalite kontrollerinin analiz komutlarından önce yer alması
önemlidir. \texttt{assert} satırları bir sorun olduğunda işlemi durdurur.
Gerçek projede veri kaynağı, paket sürümleri ve rassal tohum da kaydedilmelidir.

\section{Python ile görselleştirme ve çıktı kaydı}

\begin{lstlisting}
import matplotlib.pyplot as plt

fig, ax = plt.subplots(1, 3, figsize=(11, 3.2))

# Gruplara gore son-test noktalarini goster
for i, g in enumerate(["A", "B"]):
    y = veri.loc[veri["grup"] == g, "son"]
    ax[0].scatter([i] * len(y), y, label=g)
ax[0].set_xticks([0, 1], ["A", "B"])
ax[0].set_ylabel("Son-test puani")

# Bireysel on-son degisimleri
for _, satir in veri.iterrows():
    ax[1].plot([0, 1], [satir["on"], satir["son"]],
               marker="o", alpha=.65)
ax[1].set_xticks([0, 1], ["On", "Son"])
ax[1].set_ylabel("Basari puani")

# Sacilim ve regresyon dogrusu
ax[2].scatter(veri["saat"], veri["son"])
sira = np.linspace(veri["saat"].min(), veri["saat"].max(), 100)
ax[2].plot(sira, basit.predict(pd.DataFrame({"saat": sira})))
ax[2].set_xlabel("Haftalik calisma saati")
ax[2].set_ylabel("Son-test puani")

plt.tight_layout()
plt.savefig("laboratuvar-sekiller.png", dpi=300,
            bbox_inches="tight")
plt.show()
\end{lstlisting}

Dosya adı, çözünürlük ve grafiği üreten kod analiz kaydının parçasıdır.
Grafikler elle düzenlenip veriyle bağı koparılmamalı; mümkün olduğunca koddan
yeniden üretilebilmelidir.

\section{SPSS işlem ve doğrulama haritası}

\begin{description}
  \item[Veri denetimi] Analyze $>$ Descriptive Statistics $>$ Frequencies / Explore
  \item[Değişim puanı] Transform $>$ Compute Variable
  \item[Eşleştirilmiş test] Analyze $>$ Compare Means $>$ Paired-Samples T Test
  \item[Bağımsız test] Analyze $>$ Compare Means $>$ Independent-Samples T Test
  \item[Korelasyon] Analyze $>$ Correlate $>$ Bivariate
  \item[Regresyon] Analyze $>$ Regression $>$ Linear
  \item[Çapraz tablo] Analyze $>$ Descriptive Statistics $>$ Crosstabs
\end{description}

SPSS menüleriyle yapılan işlemler \textbf{Paste} düğmesi kullanılarak syntax
dosyasına aktarılmalıdır. Syntax, hangi değişkenin hangi alana konduğunu,
eksik veri seçeneklerini ve analiz sırasını denetlenebilir hâle getirir.

\begin{outputguidebox}
Çıktıyı test test okumak yerine araştırma sorusu sırasıyla inceleyin. Her
soruda önce geçerli $n$ ve betimsel istatistikleri, sonra tahmin ve güven
aralığını, ardından test ve etki büyüklüğünü, son olarak varsayım tanılarını
okuyun. SPSS Viewer dosyası tek başına yeniden üretilebilir analiz değildir;
syntax ve veri işleme günlüğü de saklanmalıdır.
\end{outputguidebox}

\section{Analiz günlüğü ve proje yapısı}
\index{analiz günlüğü}
\index{proje yapısı}
\index{yeniden üretilebilirlik}

Tekrarlanabilir bir proje en az şu öğeleri içerir:

\begin{itemize}
  \item değiştirilmeyen ham veri,
  \item kimlikten arındırılmış analiz verisi,
  \item veri sözlüğü,
  \item temizleme ve türetme kodu,
  \item analiz ve grafik kodu,
  \item yazılım ile paket sürümleri,
  \item alınan kararların tarihli analiz günlüğü,
  \item üretilen tablo, şekil ve rapor.
\end{itemize}

Ham veri üzerinde elle yapılan ve kaydedilmeyen değişiklikler yeniden
üretilebilirliği bozar. Temiz veri her seferinde ham veriden kodla
oluşturulmalıdır. Katılımcı kimlikleri analiz dosyasından ayrı ve güvenli
tutulmalıdır.

\section{Sonuçları tek anlatıda birleştirmek}

Ayrı analizler birbiriyle çelişiyor gibi görünebilir:

\begin{itemize}
  \item bütün öğrencilerde güçlü ortalama artış vardır,
  \item son-test grup farkı küçük örneklemde belirsizdir,
  \item değişim farkı gruplar arasında belirgindir,
  \item çalışma süresi son-testle çok güçlü ilişkilidir.
\end{itemize}

Bu sonuçlar farklı hedefleri yanıtladığı için doğrudan çelişmez. Bütün
öğrencilerin zaman içindeki değişimi, iki grubun dönem sonu konumu, grupların
artış farkı ve öğrenciler arası çalışma--başarı ilişkisi ayrı niceliklerdir.

Araştırmacı analizleri tek bir nedensel hikâyeye zorlamamalıdır. Grup ataması
ve çalışma süresi gözlemselse hem grup farkları hem regresyon eğimi karıştırıcı
değişkenlerden etkilenebilir.

\section{Bütünleştirici yöntem paragrafı}

``Her satırın bir öğrenciyi temsil ettiği veri setinde ön-test, son-test ve
haftalık çalışma süresi incelenmiştir. Değişim puanı son-test eksi ön-test
olarak hesaplanmıştır. Veri aralıkları, yinelenen kimlikler ve eksik gözlemler
analizden önce kontrol edilmiştir. Son-test puanları betimsel istatistik ve
yüzde 95 güven aralığıyla özetlenmiş; zaman içindeki değişim eşleştirilmiş
$t$ testi, gruplar arası son-test farkı Welch $t$ testiyle değerlendirilmiştir.
Çalışma süresi ile son-test ilişkisi saçılım grafiği, Pearson korelasyonu ve
basit doğrusal regresyonla incelenmiştir. Tahminler güven aralıkları ve etki
büyüklükleriyle raporlanmıştır.''

\section{Bütünleştirici bulgu paragrafı}

``Öğrencilerin son-test puanı ortalama 59,50'dir ($SS=5{,}90$, yüzde 95 GA
$[56{,}36;62{,}64]$, $n=16$). Son-test puanları ön-testten ortalama 5,38
puan daha yüksektir, $t(15)=11{,}15$, $p<0{,}001$, yüzde 95 GA
$[4{,}35;6{,}40]$, $d_z=2{,}79$. A ve B gruplarının son-test farkı 3,50
puandır; ancak tahmin belirsizdir, Welch $t(13{,}33)=1{,}20$, $p=0{,}249$,
yüzde 95 GA $[-2{,}76;9{,}76]$, $d=0{,}60$. Çalışma süresi son-test
puanıyla güçlü pozitif ilişki göstermiştir, $r(14)=0{,}981$, $p<0{,}001$.
Regresyon eğimi saat başına 3,14 puandır (yüzde 95 GA $[2{,}78;3{,}50]$,
$R^2=0{,}962$). Küçük ve gözlemsel öğretim verisi nedeniyle bulgular nedensel
etki veya dış evrene genelleme olarak yorumlanmamıştır.''

\section{Etik ve ölçülü yorum}
\index{araştırma etiği}
\index{mahremiyet}
\index{damgalayıcı dil}

PDR verileri çoğu zaman hassastır. İyi analiz yalnız doğru formül değil,
katılımcı mahremiyeti ve zarar vermeyen raporlama gerektirir:

\begin{itemize}
  \item küçük hücrelerde kimliği açığa çıkarabilecek ayrıntılar verilmemeli,
  \item risk kategorileri damgalayıcı dille sunulmamalı,
  \item grup ortalaması bireysel tanı gibi kullanılmamalı,
  \item belirsizlik ve sınırlılıklar görünür olmalı,
  \item sonuçlara göre sonradan değiştirilen analizler açıklanmalı,
  \item veri ve kod paylaşımı onam ile gizlilik sınırlarına uymalıdır.
\end{itemize}

\begin{researchbox}
Tekrarlanabilirlik veriyi herkese açmak demek değildir. Hassas veriler
paylaşılmadan kod, sentetik örnek, değişken sözlüğü ve analiz kararları açık
tutulabilir. Bilimsel şeffaflık ile katılımcı gizliliği birlikte korunmalıdır.
\end{researchbox}

\section{Son teslim ürünü}

Bu laboratuvarın ürünü yalnız bir çıktı dosyası değil, küçük bir araştırma
paketidir:

\begin{enumerate}
  \item araştırma soruları ve yöntem eşleştirme tablosu,
  \item veri sözlüğü ve kalite kontrol günlüğü,
  \item bir betimsel tablo,
  \item en az iki açıklayıcı grafik,
  \item güven aralığı, karşılaştırma ve ilişki analizleri,
  \item varsayım ve duyarlılık değerlendirmesi,
  \item yöntem, bulgu ve sınırlılık paragrafları,
  \item yeniden çalıştırılabilir Python kodu veya SPSS syntax dosyası.
\end{enumerate}

\section{Değerlendirme anahtarı}

\begin{table}[htbp]
\centering
\caption{Bütünleştirici laboratuvar için kısa değerlendirme rubriği.}
\begingroup\scriptsize
\begin{tabular}{@{}p{0.22\textwidth}p{0.30\textwidth}p{0.30\textwidth}@{}}
\toprule
Boyut & Güçlü çalışma & Geliştirilmesi gereken çalışma \\
\midrule
Soru--yöntem uyumu & Her yöntem açık hedefe bağlıdır & Testler amaç belirtilmeden seçilmiştir \\
Veri kalitesi & Kontroller ve kararlar kaydedilmiştir & Temizleme görünmez veya sonuç odaklıdır \\
İstatistiksel doğruluk & Tahmin, GA, etki ve test uyumludur & Yalnız $p$-değeri veya yanlış test vardır \\
Görselleştirme & Ham veri ve örüntü görünürdür & Yalnız süslü özet grafik vardır \\
Yorum & Tasarım ve belirsizlik sınırları korunur & Nedensellik veya kesinlik abartılır \\
Tekrarlanabilirlik & Kod/syntax ve sürümler kaydedilmiştir & İşlemler elle ve belgelenmeden yapılmıştır \\
Etik & Gizlilik ve damgalamama gözetilmiştir & Küçük hücre veya kimlik riski vardır \\
\bottomrule
\end{tabular}
\endgroup
\end{table}

\section{Literatür veri setiyle IBM SPSS bütünleştirici laboratuvarı}
\index{Student Performance veri seti!bütünleştirici analiz}
\index{SPSS!bütünleştirici analiz}
\index{SPSS!analiz günlüğü}
\index{yeniden üretilebilirlik!SPSS syntax}

Bu uygulamada önceki bölümlerde tanıtılan \emph{Student Performance}
matematik dosyası baştan sona tek bir IBM SPSS iş akışıyla analiz edilecektir
\cite{cortez2008data,cortezsilva2008}. Veri iki Portekiz okulundaki
öğrencilerin notlarını, çalışma süresi kategorilerini ve bazı demografik,
sosyal ve okul değişkenlerini içerir. Bu laboratuvarda yalnız şu değişkenler
kullanılacaktır:

\begin{description}
  \item[\textnormal{\texttt{school}}] Öğrencinin okulu: GP veya MS,
  \item[\textnormal{\texttt{studytime}}] Haftalık çalışma süresi: 1'den 4'e sıralı kategori,
  \item[\textnormal{\texttt{G3}}] 0--20 aralığındaki yıl sonu matematik notu,
  \item[\textnormal{\texttt{basarili}}] Bu laboratuvarda \texttt{G3}$\geq10$ olarak türetilen gösterge.
\end{description}

Amaç en fazla sayıda anlamlı sonuç bulmak değil, araştırma sorularının veri
denetiminden raporlamaya kadar aynı syntax içinde izlenebilmesini sağlamaktır.

\begin{researchbox}
Matematik dosyasında 395 öğrenci bulunması, örneklemin bütün Portekizli
öğrencileri temsil ettiği anlamına gelmez. Veri iki okulla sınırlıdır ve
öğrencilerin ülke genelinden olasılıklı yöntemle seçildiği belirtilmemiştir.
Bu nedenle aşağıdaki standart hatalar ve testler örneklem içi belirsizliği
gösterir; seçim yanlılığını ortadan kaldırmaz.
\end{researchbox}

\subsection{Araştırma soruları ve önceden tanımlanan analizler}

\begin{table}[htbp]
\centering
\caption{Literatür verisi için bütünleştirici analiz planı.}
\begin{tabular}{@{}p{0.43\textwidth}p{0.40\textwidth}@{}}
\toprule
Araştırma sorusu & SPSS yordamı \\
\midrule
G3 puanlarının merkezi, yayılımı ve biçimi nasıldır? & Frequencies ve Explore \\
Ortalama G3 puanı 10 referansından farklı mıdır? & Tek örneklem $t$ testi \\
G3 ortalamaları çalışma süresi gruplarında farklı mıdır? & Tek yönlü ve Welch ANOVA \\
Okul ile başarı durumu ilişkili midir? & Crosstabs ve Pearson ki-kare \\
Sonuçlar hangi etki büyüklükleriyle tamamlanmalıdır? & Cohen $d$, eta-kare ve Cramér $V$ \\
\bottomrule
\end{tabular}
\end{table}

Birincil soru veri görülmeden önce belirlenmelidir. Burada çok sayıda yöntem
aynı veri üzerinde öğretim amacıyla gösterildiği için bütün sonuçlar
keşfedici bir analiz ailesinin parçaları olarak değerlendirilecektir.

\subsection{Tek dosyada yeniden üretilebilir SPSS syntax}

Birinci bölümde içe aktarılan \texttt{StudentMath} veri seti açıkken aşağıdaki
syntax çalıştırılır. Önce çalışma kopyası oluşturulur; böylece özgün içe
aktarım dosyası değiştirilmeden korunur.

\begin{lstlisting}[language={}]
* 1. Calisma kopyasi ve veri denetimi.
DATASET ACTIVATE StudentMath.
DATASET COPY StudentMathLab.
DATASET ACTIVATE StudentMathLab.

DISPLAY DICTIONARY.

FREQUENCIES VARIABLES=school studytime
 /ORDER=ANALYSIS.

FREQUENCIES VARIABLES=G3
 /STATISTICS=MEAN MEDIAN STDDEV MINIMUM MAXIMUM QUARTILES
 /HISTOGRAM NORMAL
 /ORDER=ANALYSIS.

COMPUTE g3_aralik_hatasi=
 (NOT MISSING(G3) AND (G3<0 OR G3>20)).
VARIABLE LABELS g3_aralik_hatasi
 'G3 degeri 0-20 araligi disinda'.
VALUE LABELS g3_aralik_hatasi 0 'Uygun' 1 'Aralik disi'.
FREQUENCIES VARIABLES=g3_aralik_hatasi.

* 2. Basari gostergesini eksikleri koruyarak uret.
COMPUTE basarili=$SYSMIS.
IF NOT MISSING(G3) basarili=(G3>=10).
VARIABLE LABELS basarili 'Yil sonu notu en az 10'.
VALUE LABELS basarili 0 'Basarisiz' 1 'Basarili'.
VARIABLE LEVEL basarili (NOMINAL).
FORMATS basarili (F1.0).
EXECUTE.

* 3. Dagilim ve grup grafiklerini incele.
EXAMINE VARIABLES=G3 BY studytime
 /PLOT BOXPLOT HISTOGRAM NPPLOT
 /STATISTICS DESCRIPTIVES
 /CINTERVAL 95
 /MISSING LISTWISE.

* 4. Ortalama G3 puanini 10 referansiyla karsilastir.
T-TEST
 /TESTVAL=10
 /MISSING=ANALYSIS
 /VARIABLES=G3
 /CRITERIA=CI(.95).

* Tek orneklem Cohen d.
AGGREGATE
 /OUTFILE=* MODE=ADDVARIABLES
 /BREAK=
 /n_G3=N(G3)
 /ortalama_G3=MEAN(G3)
 /ss_G3=SD(G3).
COMPUTE d_tek=(ortalama_G3-10)/ss_G3.
FORMATS ortalama_G3 ss_G3 d_tek (F7.3).

* 5. Calisma suresi gruplarini karsilastir.
ONEWAY G3 BY studytime
 /STATISTICS DESCRIPTIVES HOMOGENEITY WELCH
 /MISSING ANALYSIS.

* 6. Okul ile basari durumunun capraz tablosu.
CROSSTABS
 /TABLES=school BY basarili
 /FORMAT=AVALUE TABLES
 /STATISTICS=CHISQ PHI RISK
 /CELLS=COUNT EXPECTED ROW COLUMN ASRESID
 /COUNT ROUND CELL.

* 7. Analiz verisini ve Viewer ciktisini kaydet.
SAVE OUTFILE='student-math-lab.sav'.
OUTPUT SAVE OUTFILE='student-math-lab.spv'.
\end{lstlisting}

Dosya adları göreli yazılmıştır; SPSS bunları etkin çalışma klasörüne
kaydeder. Kurumsal veya kişisel bilgisayarda proje klasörü önceden
belirlenmeli, ham veriyle üretilen dosyalar farklı alt klasörlerde tutulmalıdır.

\subsection{Kontrol noktaları}

Syntax sırasındaki denetimler analiz sonuçları kadar önemlidir:

\begin{enumerate}
  \item \textbf{Dictionary} tablosunda değişken türü, etiketi ve eksik değer tanımları kontrol edilir.
  \item Frekans tablolarında geçerli $n$, eksik sayısı ve kategori kodları doğrulanır.
  \item G3 için 0--20 dışında değer bulunmadığı denetlenir.
  \item \texttt{basarili} değişkenindeki 0 ve 1 sayıları özgün G3 değerleriyle karşılaştırılır.
  \item Histogram, Q--Q ve kutu grafikleri test tablolarından önce incelenir.
  \item Çıktıdaki grup sırası, farkların ve odds oranının yönüyle eşleştirilir.
\end{enumerate}

\begin{mistakebox}
Çalışma kopyası oluşturmak tek başına yeniden üretilebilirlik sağlamaz.
Syntax satırları seçilip farklı sırayla çalıştırılırsa eski türetilmiş
değişkenler veya açık kalan ağırlık ve filtreler sonuçları etkileyebilir.
Ana syntax temiz bir oturumda baştan sona çalıştırılmalı ve Viewer çıktısı bu
çalıştırmayla birlikte kaydedilmelidir.
\end{mistakebox}

\subsection{SPSS çıktılarının bütünleştirici özeti}

\begin{table}[htbp]
\centering
\caption{Student Performance analizlerinin ortak sonuç özeti.}
\begingroup\small
\begin{tabular}{@{}p{0.28\textwidth}p{0.25\textwidth}p{0.31\textwidth}@{}}
\toprule
Analiz & Nokta tahmini & Çıkarımsal sonuç \\
\midrule
G3 betimsel özeti & $\bar x=10{,}415$, $SS=4{,}581$, $Md=11$ & Yüzde 95 GA $[9{,}962;10{,}868]$ \\
Tek örneklem testi & 10'a göre fark $0{,}415$ & $t(394)=1{,}801$, $p=0{,}072$, $d=0{,}091$ \\
Çalışma süresi ANOVA & Grup ortalamaları 10,048--11,400 & $F(3,391)=1{,}728$, $p=0{,}161$, $\eta^2=0{,}013$ \\
Okul--başarı ilişkisi & GP \%67,6; MS \%63,0 başarılı & $\chi^2(1,N=395)=0{,}386$, $p=0{,}534$, $V=0{,}031$ \\
\bottomrule
\end{tabular}
\endgroup
\end{table}

G3 ortalaması 10 referansından 0,415 puan yüksektir; fakat güven aralığı sıfır
farkını içerir ve $p=0{,}072$'dir. Bu sonuç ortalamanın kesin olarak 10
olduğunu kanıtlamaz. $d=0{,}091$, gözlenen farkın standart sapmaya göre çok
küçük olduğunu gösterir.

Çalışma süresi gruplarının örneklem ortalamaları farklı görünse de grup içi
değişkenliğe göre genel fark belirgin değildir. Eta-kare yaklaşık 0,013'tür;
örneklemdeki G3 değişkenliğinin yaklaşık yüzde 1,3'ü çalışma süresi
kategorileri arasındaki ayrımla ilişkilidir. Çalışma süresinin geniş ve öz
bildirime dayalı kategorilerle ölçülmesi bu sonucun yorum sınırıdır.

GP okulunda başarı oranı yüzde 67,6, MS okulunda yüzde 63,0'dır. Pearson
ki-kare testi ve çok küçük Cramér $V$ değeri, bu örneklemde okul ile ikili
başarı göstergesi arasında belirgin bir ayrışma göstermemektedir. Bu sonuç
okulların eşdeğer olduğunu kanıtlamaz; özellikle MS grubunun daha küçük olması
tahmin hassasiyetini sınırlar.

\begin{outputguidebox}
Bütünleştirici Viewer dosyasını tablo sırasıyla değil, araştırma sorusu
sırasıyla okuyun. Her soru için önce geçerli $n$ ve dağılımı, sonra nokta
tahmini ile güven aralığını, ardından test istatistiği ve $p$-değerini, son
olarak etki büyüklüğü ile tasarım sınırlarını kaydedin. Levene, normallik veya
beklenen frekans tablolarını ana araştırma sonucu gibi sunmayın.
\end{outputguidebox}

\subsection{Bütünleştirici yöntem ve bulgu metni}

``Student Performance matematik dosyasında değişken türleri, kategori kodları,
eksik gözlemler ve G3 puan aralığı analizden önce denetlenmiştir. Yıl sonu notu
betimsel istatistik ve yüzde 95 güven aralığıyla özetlenmiş; ortalama 10
referansıyla tek örneklem $t$ testi kullanılarak karşılaştırılmıştır. Çalışma
süresi grupları klasik ve Welch ANOVA, okul ile önceden tanımlanan başarı
durumu ise çapraz tablo ve Pearson ki-kare testiyle incelenmiştir. Sonuçlar
Cohen $d$, eta-kare ve Cramér $V$ ile tamamlanmıştır.''

``Yıl sonu matematik notu ortalama 10,42'dir ($SS=4{,}58$, yüzde 95 GA
$[9{,}96;10{,}87]$, $n=395$) ve 10 referansından belirgin biçimde farklı
değildir, $t(394)=1{,}80$, $p=0{,}072$, $d=0{,}09$. Çalışma süresi
gruplarının ortalamaları için yeterli fark kanıtı bulunmamıştır,
$F(3,391)=1{,}73$, $p=0{,}161$, $\eta^2=0{,}013$. Benzer biçimde okul ile
başarı durumu arasındaki ilişki belirgin değildir,
$\chi^2(1,N=395)=0{,}386$, $p=0{,}534$, $V=0{,}031$. Analizler keşfedici
olarak değerlendirilmiş ve iki okulla sınırlı veri ülke geneline
genellenmemiştir.''

Bu metin farklı analizleri tek bir ``etki var veya yok'' sonucuna indirgemez.
Her analiz farklı hedef büyüklüğü yanıtlar ve anlamlı olmayan sonuçlar eşitlik
kanıtı olarak kullanılmaz.

\section{R ile bütünleştirici analiz akışı}
\label{sec:r-integrated-lab}

\textbf{Soru.} Küçük öğretim verisinde veri denetimi, betimsel özet,
eşleştirilmiş değişim, bağımsız grup karşılaştırması, korelasyon ve başlangıç
puanı ayarlı modeli tek bir R akışında üretin.

\begin{lstlisting}[style=prismR]
veri <- data.frame(
  id = 1:12,
  grup = factor(rep(c("A", "B"), each = 6)),
  on = c(40,45,50,55,60,65, 42,47,52,57,62,67),
  son = c(46,49,55,62,63,71, 43,47,54,58,61,67)
)
stopifnot(anyDuplicated(veri$id) == 0)
stopifnot(all(complete.cases(veri)))
veri$degisim <- veri$son - veri$on

summary(veri)
aggregate(cbind(on, son, degisim) ~ grup,
          data = veri, FUN = mean)

t.test(veri$son, veri$on, paired = TRUE)
t.test(son ~ grup, data = veri, var.equal = FALSE)
cor.test(veri$on, veri$son)

ayarlanmis <- lm(son ~ on + grup, data = veri)
summary(ayarlanmis)
confint(ayarlanmis)
plot(ayarlanmis)

write.csv(veri, "analiz-verisi.csv", row.names = FALSE)
capture.output(sessionInfo(), file = "session-info.txt")
\end{lstlisting}

\begin{outputguidebox}
\textbf{Çözüm ve kontrol değerleri.} Ortalama değişim 2,833 puandır;
eşleştirilmiş test $t(11)=3{,}602$, $p=0{,}00415$ verir. Son-test için
Welch grup karşılaştırması $t\approx0{,}504$, $p\approx0{,}626$'dır.
Ön-test ile son-test korelasyonu $r=0{,}953$'tür. Başlangıç ayarlı modelde
grup katsayısının işareti, referans grubun düzeyine bağlıdır; katsayı adı ve
faktör kodlaması görülmeden yorum yapılmamalıdır.
\end{outputguidebox}

\textbf{Yorum.} Eşleştirilmiş değişim, son-test grup farkı ve başlangıç ayarlı
model aynı soruyu yanıtlamaz. Kodun tek dosyada çalışması yöntem seçimini
otomatik olarak doğrulamaz. Veri sözlüğü, dışlama kuralları, eksik veri kararı
ve analiz planı kodla birlikte saklanmalıdır.

\textbf{Pekiştirme ve yanıt.} Grup etiketlerinin sırasını değiştirerek modeli
yeniden kurun. Uydurulan değerler değişmez; kesme ve grup katsayısının
parametreleştirilmesi değişir. Raporlamada hangi grubun referans olduğu açıkça
belirtilmelidir.
\section{Bölüm sonu çözümlü problemler}

\begin{enumerate}
  \item \textbf{Soru--yöntem eşleştirmesi.} Aynı öğrencilerin ön-test ve son-test puanları, iki bağımsız program grubunun son-test puanları ve çalışma saatiyle başarı puanı incelenecektir. Üç temel yöntemi seçin.

  \textbf{Çözüm:} Aynı öğrencilerin değişimi için eşleştirilmiş $t$ testi; bağımsız grupların ortalama farkı için Welch $t$ testi; iki nicel değişkenin doğrusal ilişkisi için saçılım grafiğiyle birlikte Pearson korelasyonu veya basit doğrusal regresyon kullanılır.

  \item \textbf{Tutarsız görünen sonuçlar.} Grup farkı için $d=0{,}60$, yüzde 95 güven aralığı $[-0{,}2;1{,}4]$ ve $p=0{,}14$ bulunmuştur. Sonucu bütünleştirin.

  \textbf{Çözüm:} Nokta tahmini orta büyüklükte görünse de aralık küçük ters etkiden büyük pozitif etkiye kadar geniştir ve sıfırı içerir. $p$-değeri seçilen eşikte kanıtın sınırlı olduğunu gösterir. Sonuç ``etki yoktur'' değil, tahmin hassasiyetinin yetersiz olduğu biçiminde sunulmalıdır.

  \item \textbf{Tekrarlanabilir analiz paketi.} Ham veri üzerinde iki elle düzeltme yapılmış, fakat değişiklikler kaydedilmemiştir. Süreci yeniden üretilebilir hâle getirmek için ne yapılmalıdır?

  \textbf{Çözüm:} Ham dosya değişmeden korunmalı; düzeltmeler doğrulama kuralları ve gerekçeleriyle kodlanarak temiz veri yeniden üretilmelidir. Veri sözlüğü, analiz günlüğü, yazılım sürümleri, kod, rassal tohumlar ve çıktı üretim adımları aynı proje yapısında saklanmalıdır.
\end{enumerate}

\bolumkaynaklari{\cite{mccullagh2022,james2021,anscombe1973,efron1993,cortez2008data,cortezsilva2008}}

\section*{Hatırla--Uygula--Yansıt}

\begin{enumerate}
\renewcommand{\labelenumi}{\textbf{\Alph{enumi}.}}
  \item \textbf{Hatırla:} Soru--tasarım--veri--yöntem--yorum zincirindeki her aşamanın görevini açıklayın.
  \item \textbf{Denetle:} Veri seti için en az sekiz otomatik veya görsel kalite kontrolü yazın.
  \item \textbf{Eşleştir:} Yedi araştırma sorusunun her birini hedef büyüklük ve uygun yöntemle eşleştirin.
  \item \textbf{Betimle:} Son-test puanlarını merkez, yayılım, biçim ve olağandışı gözlemler açısından özetleyin.
  \item \textbf{Aralık:} Son-test ortalamasının güven aralığını hesaplayıp doğru ve yanlış birer yorum yazın.
  \item \textbf{Eşleştirilmiş analiz:} Farkın yönünü, tahminini, güven aralığını ve $d_z$ değerini yorumlayın.
  \item \textbf{Bağımsız analiz:} $d=0{,}60$ iken $p=0{,}249$ bulunmasının neden çelişki olmadığını açıklayın.
  \item \textbf{Yöntem karşılaştır:} Son-test, değişim puanı ve başlangıç ayarlı modelin farklı hedeflerini ayırın.
  \item \textbf{İlişki:} Çok yüksek korelasyonun hangi veri ve tasarım kontrollerini gerekli kıldığını yazın.
  \item \textbf{Regresyon:} Eğim, güven aralığı, $R^2$ ve artık hatasını ayrı ayrı yorumlayın.
  \item \textbf{Kategorileştirme:} Değişim puanını eşikle ikiye ayırmanın bilgi ve güç kaybını açıklayın.
  \item \textbf{Eksik veri:} B grubundaki düşük puanlı öğrencilerin kaybının grup karşılaştırmasını nasıl yanlılaştırabileceğini tartışın.
  \item \textbf{Duyarlılık:} En az iki ana analiz için alternatif yöntem ve hangi sorunu değerlendirdiğini yazın.
  \item \textbf{Uygula:} Python kodunu çalıştırıp tablo ve grafiklerin metindeki sonuçlarla tutarlılığını denetleyin.
  \item \textbf{Tekrarla:} Analizi SPSS syntax ile yeniden üretip Python sonuçlarıyla karşılaştırın.
  \item \textbf{Raporla:} Yöntem, bulgu, sınırlılık ve etik paragrafından oluşan kısa bir araştırma özeti hazırlayın.
  \item \textbf{Yansıt:} Analiz sürecinde en fazla yargı gerektiren üç kararı ve farklı bir araştırmacının neden başka karar verebileceğini açıklayın.
\end{enumerate}